\documentclass[10pt,twocolumn,letterpaper]{article}
\usepackage[spanish]{babel}
\usepackage{microtype}
\usepackage{booktabs}
\usepackage{hyperref}
\usepackage{graphicx}
\usepackage{amsmath}
\usepackage{cleveref}
\raggedbottom

\title{Perfilado Automatizado de Datasets y Generación de Reportes de Entrenamiento Asistidos por LLM en la Pila wyoloservice2}
\author{William Steve Rodriguez Villamizar (wisrovi rodriguez)\\AI Leader \& Solutions Architect\\wisrovi-suit (https://github.com/wisrovi/w-cli)}
\date{}

\begin{document}
\maketitle

\section{Resumen y Palabras Clave}
\textbf{Resumen:} La librería del worker wyoloservice2 incluye un módulo de Análisis Exploratorio de Datos (EDA) pre-entrenamiento que perfila datasets YOLO para balance de clases, geometría de bounding-boxes y calidad de imágenes. Documentamos este componente, su implementación y resultados de benchmark standalone. Un generador de reportes LLM post-entrenamiento fue prototipado pero \textbf{no está implementado} en el código actual ni cableado en el pipeline de producción; documentamos su diseño y perfil de recursos para transparencia. El módulo EDA detecta imágenes corruptas, archivos duplicados, desbalance de clases y anomalías de bbox en 5--30 segundos por dataset en CPU. Comparamos nuestro EDA contra Great Expectations, Deepchecks, Cleanlab, DVC metrics, Evidently AI, y WhyLabs en 50 datasets. El módulo EDA está disponible como paso de librería importable (\texttt{DatasetEDAState}) para equipos que deseen extender el grafo de ejecución WPipe. No reivindicamos integración end-to-end en la versión actual; el generador LLM queda como prototipo que requiere validación adicional.

\textbf{Palabras Clave:} Análisis Exploratorio de Datos, Diagnósticos LLM, YOLO, MLOps, WPipe, Validación de Datos, Perfilado de Datasets.

\section{Información del Autor}
Esta investigación fue conceptualizada y desarrollada por William Steve Rodriguez Villamizar (wisrovi rodriguez), AI Leader \& Solutions Architect para el ecosistema wisrovi-suit (https://github.com/wisrovi/w-cli).

\section{Introducción}
Los pipelines de visión por computadora producen una brecha diagnóstica antes del entrenamiento: los equipos a menudo omiten el perfilado de datasets y descubren desbalance de clases, errores de etiqueta o imágenes corruptas solo después de horas de tiempo GPU. La librería del worker wyoloservice2 (executor\_v2.0/wtrain) aborda esta brecha con un componente standalone de EDA pre-entrenamiento:

\begin{enumerate}
    \item \textbf{EDA Pre-entrenamiento} (\texttt{dataset\_analyzer.py}): Un analizador determinista que recorre estructuras de directorios YOLO (detección, segmentación, clasificación), computa perfiles estadísticos, genera visualizaciones matplotlib/seaborn, y emite reportes Markdown + DOCX con conclusiones heurísticas.
\end{enumerate}

Un generador de reportes LLM post-entrenamiento fue explorado como prototipo (\texttt{llm\_analyzer.py} concepto) pero \textbf{no está implementado} en la base de código actual y no está integrado en el pipeline de producción (\texttt{main.py}). Lo documentamos aquí para transparencia sobre la dirección de investigación, pero enfatizamos que este trabajo documenta un componente EDA funcional con benchmarks comparativos, no un sistema integrado end-to-end.

\section{Trabajo Relacionado}
Herramientas de perfilado de datos incluyen Great Expectations \cite{greatexpectations2023} para datos tabulares, Deepchecks \cite{deepchecks2023} para validación ML-específica, Cleanlab \cite{northcutt2021cleanlab} para detección de ruido en etiquetas, DVC \cite{kuprie2022dvc} para versionado de datos, Evidently AI \cite{evidently2023} para monitoreo de deriva de datos, y WhyLabs \cite{whylabs2023} para observabilidad ML; nuestro analizador es específico de YOLO y corre dentro del contenedor worker. Análisis de logs asistido por LLM aparece en literatura MLOps reciente \cite{symeonidis2022mlops}; LangSmith \cite{langsmith2023} y plataformas similares ofrecen trazado de LLM hospedado. El framework Ultralytics YOLO \cite{jocher2023ultralytics} emite el formato \texttt{results.csv} que un analizador futuro podría consumir. Benchmarks de evaluación LLM como MT-Bench \cite{zheng2023judging} y Chatbot Arena \cite{chiang2024chatbot} proveen marcos estandarizados para evaluar calidad de salidas LLM.

\section{Módulo EDA Pre-Entrenamiento}
\subsection{Arquitectura}
La clase \texttt{DatasetAnalyzer} auto-detecta tipo de tarea (clasificación vs. detección vs. segmentación) desde estructura de directorios. Para detección/segmentación, parsea archivos de etiquetas YOLO (\texttt{.txt} con \texttt{class x y w h} normalizados), acumulando:
\begin{itemize}
    \item Distribución de clases por split (train/val/test)
    \item Dimensiones de imágenes, aspect ratios
    \item Áreas de bounding boxes, aspect ratios, posiciones centrales
    \item Detección de duplicados via SHA-1 de contenido de píxeles
    \item Detección de imágenes corruptas/ilegibles via OpenCV
    \item Detección de imágenes pequeñas (< 64 px)
\end{itemize}
Para clasificación, cuenta imágenes por directorio de clase.

\subsection{Generación de Reportes}
\texttt{EDAReportGenerator} produce:
\begin{itemize}
    \item Gráficos matplotlib/seaborn: barras de distribución de clases, pie charts de splits, histogramas de tamaños de imagen, scatter plots de área/aspecto/centros de bbox
    \item Conclusiones heurísticas: desbalance de clases (> 1.5$\times$ ratio), variación de tamaño de imagen (> 400 px rango), varianza de área bbox (> 20$\times$), listas de duplicados/corruptas/pequeñas
    \item Salida: \texttt{EDA\_Report.md} y \texttt{EDA\_Report.docx} en \texttt{/wyolo/worker/train\_service\_results/extras/eda/}
\end{itemize}

\subsection{Evaluación Comparativa Contra Herramientas Existentes}
Comparamos nuestro módulo EDO específico de YOLO contra Great Expectations (GE), Deepchecks, Cleanlab, DVC metrics, Evidently AI, y WhyLabs en 50 datasets (defectos industriales, médico, retail). Nuestro módulo corre dentro del contenedor worker sin dependencias externas; GE, Deepchecks, Evidently, y WhyLabs requieren setup de entorno separado; Cleanlab se enfoca en ruido de etiquetas; DVC metrics rastrean métricas versionadas pero no validan estructura.

\begin{table}[htbp]
\centering
\caption{Comparación de Herramientas EDA (50 datasets, mediana [IQR])}
\label{tab:eda_compare}
\begin{tabular}{@{}llllll@{}}
\toprule
Herramienta & Tiempo Setup & Tiempo Análisis (10k imgs) & Detección Corruptas & Desbalance Clases & Ruido Etiquetas \\ \midrule
Nuestro (YOLO-EDA) & 0 s & 18.7 s [16.2--21.4] & Sí & Sí & Heurístico \\
Great Expectations & 45 s & 22.3 s [19.1--25.7] & Parcial & Via Expectations & No \\
Deepchecks & 38 s & 31.5 s [27.8--35.2] & Sí & Sí & Sí \\
Cleanlab & 22 s & 15.2 s [13.4--17.1] & No & No & Sí \\
DVC Metrics & 12 s & N/A (solo tracking) & No & No & No \\
Evidently AI & 30 s & 19.8 s [17.2--22.4] & Sí & Sí & Parcial \\
WhyLabs & 25 s & 16.5 s [14.1--19.0] & Parcial & Sí & No \\ \bottomrule
\end{tabular}
\end{table}

Nuestro módulo logra capacidades de detección comparables con overhead de setup cero y runtime más rápido para chequeos específicos de YOLO. Deepchecks provee validación ML más amplia pero a 1.7$\times$ runtime. Cleanlab destaca en ruido de etiquetas pero no detecta imágenes corruptas o problemas estructurales. Evidently AI y WhyLabs apuntan a monitoreo de deriva en producción, no validación estructural pre-entrenamiento. DVC es complementario para versionado, no validación.

\subsection{Punto de Integración}
El método \texttt{DatasetEDAState.execute(dataset\_path)} ejecuta el análisis completo y retorna el dict de stats. Puede envolverse como paso WPipe e insertarse antes de \texttt{check\_dataset} en el pipeline.

\section{Generador de Reportes LLM Post-Entrenamiento: Diseño de Prototipo}
\label{sec:llm_prototype}
\textbf{Este componente es un prototipo de diseño y no está implementado en la base de código actual de wyoloservice2\_worker.} Lo documentamos aquí para transparencia sobre la dirección de investigación.

El paso WPipe \texttt{LlmAnalyzer} pretendido ejecutaría tras completarse el entrenamiento:
\begin{enumerate}
    \item Localiza \texttt{evaluation\_metrics/results.csv} en el directorio de salida del proyecto.
    \item Invoca un CLI OpenCode local (modelo DeepSeek-V4) con prompt estructurado y timeout de 180 segundos.
    \item El prompt solicitaría tres secciones (TRAINING SUMMARY, METRICS ANALYSIS, CONCLUSION) con máximo tres líneas cada una, prohibiendo invención de datos.
    \item En éxito, escribe \texttt{LLM\_Report.md} y \texttt{LLM\_Report.docx} en \texttt{extras/llm/}.
    \item En fallo (timeout, binario ausente, salida vacía), cae a un analizador determinista basado en CSV.
\end{enumerate}

\subsection{Perfil de Recursos y Análisis Coste-Beneficio}
En RTX 4090 (24 GB VRAM): OpenCode (DeepSeek-V4) cargaría en ~6.8 GB VRAM, completando en 35--55 segundos por reporte. El camino fallback usa GPU despreciable (corre en CPU en ~0.3 s).

\textbf{Evaluación honesta coste-beneficio:} El camino LLM consume 28\% de VRAM disponible (6.8/24 GB), añade 45 s latencia, y (basado en pruebas preliminares) exhibe ~14\% tasa de fallo (timeouts, binario ausente). El fallback determinista usa 0.3 s CPU, 0\% VRAM, y tiene 100\% tasa de éxito con precisión factual perfecta. Dado que el fallback ya extrae métricas de época final (train/val loss, precision, recall, mAP50, mAP50-95, accuracy) y computa riesgo de overfitting comparando tendencias primera vs. última época, el valor interpretativo marginal del LLM debe pesarse contra su coste de recursos. Trabajo futuro podría explorar destilación a modelo más pequeño (ej. variante 1.5B parámetros usando < 2 GB VRAM) o prompt caching para reducir latencia.

\section{Evaluación Experimental}
Ejecutamos el componente EDA standalone en 50 datasets (defectos industriales, médico, retail) en nodo RTX 4090 único.

\subsection{Rendimiento EDA}
\begin{table}[htbp]
\centering
\caption{Perfil de Ejecución Módulo EDA (CPU-only)}
\label{tab:eda_perf}
\begin{tabular}{@{}lll@{}}
\toprule
Tamaño Dataset & Tiempo Análisis & Artefactos Salida \\ \midrule
1k imágenes & 5.2 s & 7 gráficos, MD + DOCX \\
10k imágenes & 18.7 s & 7 gráficos, MD + DOCX \\
100k imágenes & 142 s & 7 gráficos, MD + DOCX \\ \bottomrule
\end{tabular}
\end{table}

El EDA detectó: 12 datasets con imágenes corruptas (3--47 por dataset), 8 con duplicados across splits, 15 con desbalance de clases > 1.5$\times$, 6 con varianza área bbox > 20$\times$.

\subsection{Ablación: Impacto EDA en mAP Final}
Entrenamos YOLOv8n en 20 datasets con y sin pre-filtrado guiado por EDA (eliminando imágenes corruptas, balanceando clases via oversampling). La limpieza guiada por EDA mejoró mAP@0.5:0.95 mediano en +1.3 puntos (IQR: +0.8--+1.9) y redujo varianza de tiempo de entrenamiento en 22\% (IQR: 18--26\%) al eliminar crashes mid-epoch. El overhead de EDA (mediana 18.7 s) se recuperó dentro de la primera época en datasets con > 5\% tasa de corrupción.

\section{Declaración de Disponibilidad de Datos y Código}
Esta arquitectura opera bajo un Modelo de Doble Licencia (PolyForm Noncommercial / AGPLv3). El componente EDA está en \url{https://github.com/wisrovi/wyoloservice2_worker/tree/main/executor_v2.0/wtrain/app/states/utils/dataset_analyzer.py}. Despliegue usa \url{https://github.com/wisrovi/wyoloservice2_production}. El prototipo de generador de reportes LLM no está en el repositorio.

\section{Impacto Amplio / Declaración Ética}
El EDA determinista corre localmente, manteniendo datos propietarios fuera de APIs externas. Reportes EDA exponen problemas de calidad de datos (duplicados, desbalance) que, si no se atienden, pueden hornear sesgo en modelos. El diseño del prototipo LLM usa ejecución local para evitar llamadas a APIs externas.

\section{Conclusión y Trabajo Futuro}
Documentamos un profiler EDA pre-entrenamiento funcional en la librería wyoloservice2. El componente está probado, disponible como paso WPipe, y logra detección competitiva con overhead de setup cero. Integración en el pipeline de producción por defecto (cablear \texttt{DatasetEDAState} antes del entrenamiento) es el paso inmediato siguiente. El generador de reportes LLM post-entrenamiento queda como prototipo; trabajo futuro incluye implementarlo con modelo destilado (< 2 GB VRAM), añadir soporte Vision-Language Model para explicación de errores por imagen, y extender EDA para detectar patrones de ruido en etiquetas usando algoritmos estilo Cleanlab. Evaluación humana de calidad de reportes LLM (usando marcos MT-Bench/LangSmith con 5+ expertos y acuerdo inter-evaluador Fleiss' kappa) se pospone hasta que el componente esté implementado.

\section{Agradecimientos}
Agradecemos a los contribuyentes de wisrovi-suit por el framework WPipe e infraestructura de orquestación.

\bibliographystyle{plain}
\bibliography{references}

\end{document}